\pagenumbering{arabic}\setcounter{page}{1}
\chapter{PENDAHULUAN}

\section{Latar Belakang}

Model bahasa besar (\textit{large language model}, LLM) semakin sering
dijalankan bukan sebagai satu model tunggal, melainkan sebagai sistem
\textit{multi-agent}: beberapa salinan model diberi peran berbeda dan bekerja
berurutan, saling meneruskan hasil pekerjaan hingga jawaban akhir terbentuk.
Pembagian peran semacam itu terbukti memperbaiki mutu jawaban pada tugas yang
menuntut perencanaan, pemeriksaan ulang, dan perbaikan bertahap.

Pada hampir seluruh sistem tersebut, medium yang menghubungkan antar-agen
adalah teks. Setiap agen harus mengubah penalaran internalnya menjadi untaian
kata, lalu agen berikutnya menguraikan kembali untaian itu menjadi representasi
internalnya sendiri. Proses bolak-balik tersebut menimbulkan dua persoalan.
Pertama, biaya komputasi meningkat karena setiap penyerahan pekerjaan
mengharuskan pembangkitan token yang panjang. Kedua, terdapat penyempitan
kanal: representasi internal yang kaya harus dimampatkan menjadi barisan token
diskret, dan hanya sebagian yang dapat dipulihkan oleh agen penerima
\citep{latentmas, cache2cache}.

Sebagai jawaban atas kedua persoalan itu, \citet{latentmas} mengusulkan
LatentMAS, kerangka kolaborasi \textit{multi-agent} yang berlangsung sepenuhnya
di ruang laten. Alih-alih bertukar teks, setiap agen menjalankan sejumlah
langkah penalaran di ruang representasi kontinu, lalu menyerahkan memori
kerjanya berupa \textit{key-value cache} dari mekanisme \textit{attention}
\citep{vaswani} kepada agen berikutnya. Pendekatan ini merupakan bagian dari
paradigma penalaran laten yang lebih luas, yang memindahkan langkah penalaran
dari barisan token diskret ke ruang representasi kontinu \citep{coconut,
latentsurvey}. Pada sejumlah tolok ukur penalaran umum, kerangka tersebut
dilaporkan menurunkan pemakaian token keluaran sebesar $70{,}8\%$--$83{,}7\%$
dan mempercepat inferensi hingga sekitar empat kali lipat, sekaligus
mempertahankan atau memperbaiki akurasi.

Seluruh mekanisme tersebut bertumpu pada satu persamaan kecil. Karena langkah
penalaran laten tidak menghasilkan token, keluaran satu langkah --- yaitu
\textit{hidden state} lapisan terakhir --- harus dipetakan kembali menjadi
sesuatu yang dapat dibaca model sebagai masukan. LatentMAS menyelesaikannya
dengan sebuah pemetaan linier yang diperoleh dari regresi \textit{ridge}
\citep{ridge} atas kedua matriks embedding model. Pemetaan itu diperkenalkan
sebagai penyelesaian rekayasa terhadap ketidaksesuaian ruang masukan dan
keluaran, dan tidak pernah dibandingkan terhadap alternatif mana pun. Dengan
kata lain, komponen yang menanggung seluruh beban mekanisme kolaborasi laten
justru merupakan komponen yang paling sedikit diuji.

Pada saat yang sama, literatur penalaran laten untuk model tunggal telah
menghasilkan beberapa cara lain membentuk masukan langkah berikutnya. Alih-alih
memetakan \textit{hidden state} secara linier, ketiganya terlebih dahulu
mengubahnya menjadi distribusi peluang atas kosakata, lalu membentuk masukan
sebagai gabungan berbobot embedding token. \textit{Soft Thinking} memakai
distribusi peluang itu secara langsung; \textit{Stochastic Soft Thinking}
\citep{softthinking} menambahkan derau Gumbel agar penalaran tidak terkunci
pada token berpeluang tertinggi; sedangkan \textit{Mixture of Inputs}
\citep{moi} menggabungkan distribusi tersebut dengan token hasil pengambilan
sampel melalui pembobotan yang bergantung pada entropi. Ketiganya dirancang dan
diuji untuk model tunggal yang bernalar sendiri, dan belum pernah dipindahkan
ke konteks kolaborasi laten antar-agen.

Terdapat alasan teoretis untuk menduga bahwa pilihan persamaan tersebut
berpengaruh, dan bahwa pengaruhnya tidak seragam antar jenis tugas. Ketiga cara
alternatif di atas menghasilkan vektor yang selalu berada di dalam lambung
konveks embedding token, yaitu wilayah yang memang pernah dilihat model selama
pelatihan. Pemetaan linier \textit{ridge} tidak memiliki kendala tersebut:
hasilnya dapat mendarat di mana saja dalam ruang berdimensi sama. Apabila
vektor yang diumpankan berada jauh dari wilayah embedding yang sah, identitas
token yang dibawanya dapat rusak. Kerusakan semacam itu masih dapat dipulihkan
pada tugas yang jawabannya pendek, tetapi tidak pada tugas yang menuntut setiap
simbol ditulis persis, seperti pembangkitan kode atau penulisan ekspresi
matematis.

Penelitian ini merumuskan kelima cara pembentukan langkah laten tersebut ke
dalam satu kerangka matematis yang sama, sehingga keempatnya dapat dinyatakan
sebagai pemilihan bobot pada simpleks peluang dan yang kelima dapat ditunjukkan
berada di luar keluarga itu. Kelimanya kemudian dibandingkan secara terkendali
pada kolaborasi \textit{multi-agent} sekuensial, dengan seluruh komponen lain
--- topologi agen, \textit{prompt}, model, dan soal --- dijaga tetap. Pengujian
dilakukan pada tiga tolok ukur penalaran umum yang mewakili tingkat tuntutan
presisi simbolik yang berbeda, ditambah satu domain simbolik berupa
pembangkitan ekspresi faktor alpha \citep{quantaalpha} sebagai pengujian
eksternal. Seluruh eksperimen dijalankan pada model lokal Qwen3-8B
\citep{qwen3}.

\section{Rumusan Masalah}

Berdasarkan latar belakang dan permasalahan yang telah diuraikan, penelitian ini dirumuskan ke dalam beberapa pertanyaan sebagai berikut.

\begin{enumerate}

      \item Bagaimana berbagai formulasi pembentukan representasi laten dapat
            dinyatakan dalam suatu kerangka matematis yang seragam sehingga
            memungkinkan perbandingan terkendali antar-metode pada sistem
            multi-agen?

      \item Bagaimana pengaruh formulasi representasi laten terhadap kemampuan
            penalaran dan fidelitas informasi simbolik pada kolaborasi multi-agen,
            dan apakah pengaruh tersebut berubah sesuai dengan tuntutan presisi
            simbolik tugas?

      \item Apakah komunikasi antar-agen melalui KV-cache mampu mempertahankan
            kinerja penalaran yang sebanding dengan komunikasi berbasis teks
            dengan biaya token dan waktu inferensi yang lebih rendah, serta
            bagaimana pengaruh formulasi representasi laten terhadap efisiensi
            tersebut?

      \item Sejauh mana representasi laten mampu mempertahankan informasi
            simbolik ketika digunakan untuk menghasilkan ekspresi faktor
            terstruktur, dan apakah kualitas keluaran yang dihasilkan tetap dapat
            dievaluasi melalui pengujian terhadap data pasar?

\end{enumerate}


\section{Pembatasan Masalah}

Agar penelitian tetap terarah dan tidak meluas di luar ruang lingkup yang
ditetapkan, penelitian ini dibatasi pada hal-hal berikut.

\begin{enumerate}

\item Model bahasa yang digunakan adalah Qwen3-8B dengan inferensi lokal.
      Penelitian tidak menguji generalisasi terhadap ukuran maupun keluarga
      model bahasa lainnya.

\item Seluruh intervensi terhadap proses penalaran bersifat
      \textit{training-free}. Tidak dilakukan perubahan bobot model,
      \textit{fine-tuning}, maupun pembelajaran penguatan.

\item Sistem multi-agen menggunakan topologi rantai sekuensial dengan
      jumlah agen yang telah ditetapkan. Topologi hierarkis, bercabang, dan
      bentuk koordinasi multi-agen lainnya tidak dievaluasi.

\item Formulasi representasi laten yang dibandingkan dibatasi pada
      \textit{raw}, \textit{soft}, \textit{sample}, \textit{Gumbel-Softmax},
      dan \textit{Mixture of Inputs} (MoI).

\item Medium komunikasi utama yang dibandingkan adalah komunikasi berbasis
      teks dan komunikasi berbasis KV-cache, dengan konfigurasi agen tunggal
      digunakan sebagai acuan. Medium gabungan KV-cache dan teks digunakan
      secara terbatas sebagai kondisi pelengkap dan tidak menjadi dasar
      utama pengujian statistik.

\item Evaluasi kemampuan penalaran pada lengan utama dibatasi pada GSM8K,
      ARC-Challenge, dan HumanEval-Plus. Ketiga tugas digunakan untuk
      merepresentasikan tuntutan keluaran numerik, pilihan simbolik, dan
      pembangkitan kode.

\item Evaluasi simbolik tambahan menggunakan generasi ekspresi faktor
      terstruktur berbasis domain-specific language (DSL). Lengan ini
      digunakan sebagai pengujian kemampuan mempertahankan muatan simbolik
      dan bukan sebagai penelitian mengenai strategi investasi.

\item Evaluasi faktor dilakukan melalui validitas ekspresi dan metrik
      kinerja faktor pada data pasar. Hasil \textit{backtest} tidak
      ditafsirkan sebagai prediksi keuntungan investasi aktual.

\end{enumerate}


\section{Tujuan dan Manfaat Penelitian}

Sejalan dengan rumusan masalah yang telah ditetapkan, penelitian ini bertujuan untuk:

\begin{enumerate}

      \item Merumuskan berbagai pendekatan pembentukan representasi laten ke dalam
            suatu kerangka matematis yang seragam sehingga karakteristik dan
            perbedaan antar-metode dapat dibandingkan secara terkendali.

      \item Menganalisis pengaruh formulasi representasi laten terhadap kemampuan
            penalaran pada sistem multi-agen serta mengidentifikasi hubungan
            antara formulasi tersebut dan kemampuan mempertahankan informasi
            simbolik pada tugas dengan tingkat tuntutan presisi yang berbeda.

      \item Mengevaluasi efisiensi komunikasi antar-agen berbasis KV-cache
            dibandingkan dengan komunikasi berbasis teks berdasarkan kinerja
            penalaran, penggunaan token, dan waktu inferensi, serta menentukan
            kondisi formulasi representasi laten yang memungkinkan komunikasi
            laten mencapai kinerja yang sebanding dengan komunikasi teks.

      \item Mengevaluasi kemampuan representasi laten dalam mempertahankan muatan
            simbolik pada tugas generasi ekspresi faktor terstruktur serta
            membandingkan keandalan dan kualitas ekspresi yang dihasilkan melalui
            evaluasi validitas ekspresi dan pengujian pada data pasar.

\end{enumerate}


Manfaat yang diharapkan adalah sebagai berikut.

\begin{enumerate}
    \item Menyediakan kerangka matematis yang menyatukan beberapa pendekatan
          penalaran laten yang selama ini dibahas terpisah, sehingga
          keduanya dapat diperbandingkan tanpa mengubah sistem di sekitarnya.
    \item Memberikan bukti empiris mengenai kondisi tempat kolaborasi laten
          dapat dan tidak dapat diandalkan, khususnya pada keluaran yang
          menuntut ketepatan simbolik.
    \item Menyediakan acuan kinerja dan biaya bagi penerapan kolaborasi laten
          pada model berukuran menengah yang dijalankan secara lokal.
\end{enumerate}

\section{Tinjauan Pustaka}

Berikut beberapa literatur kunci yang menjadi acuan penelitian ini.

\citet{latentmas} mengusulkan LatentMAS, kerangka kolaborasi
\textit{multi-agent} di ruang laten. Setiap agen membangkitkan pikiran laten
melalui representasi lapisan terakhir, lalu meneruskan memori kerjanya berupa
KV-\textit{cache} kepada agen berikutnya tanpa melalui teks. Penelitian ini
mengambil mekanisme tersebut sebagai objek yang diuji, khususnya persamaan
pemetaan \textit{hidden state} yang menjadi tumpuannya.

\citet{softthinking} menelaah mekanisme kerja \textit{soft thinking} dan
menunjukkan bahwa penalaran lunak yang deterministik cenderung terkunci pada
token berpeluang tertinggi. Sebagai perbaikan, penulis mengusulkan penambahan
derau Gumbel pada logit sebelum pembentukan bobot. Penelitian ini mengadopsi
persamaan tersebut sebagai salah satu kondisi yang dibandingkan.

\citet{moi} mengusulkan \textit{Mixture of Inputs}, yang memperlakukan
distribusi keluaran model sebagai \textit{prior} dan token hasil pengambilan
sampel sebagai observasi, lalu membentuk masukan berikutnya dari ekspektasi
\textit{posterior} keduanya. Bobot campurannya bergantung pada entropi
distribusi, sehingga menyesuaikan diri terhadap tingkat keyakinan model.

\citet{coconut} dan survei \citet{latentsurvey} menempatkan ketiga pendekatan
di atas dalam paradigma penalaran laten yang lebih luas, yaitu pemindahan
langkah penalaran dari barisan token diskret ke ruang representasi kontinu.

\citet{quantaalpha} mengusulkan kerangka penemuan faktor alpha berbasis LLM
yang memformulasikan pencarian faktor sebagai pembangkitan ekspresi simbolik
dan mengevaluasinya melalui metrik seperti \textit{information coefficient} dan
kinerja portofolio. Penelitian ini memakai domain tersebut sebagai pengujian
eksternal terhadap kemampuan kanal laten mempertahankan muatan simbolik, bukan
sebagai upaya mengembangkan kembali kerangka tersebut.

\citet{ridge} memperkenalkan regresi \textit{ridge} sebagai penyelesaian
teregularisasi bagi sistem persamaan yang berkondisi buruk. Bentuk tertutup
penyelesaian tersebut merupakan dasar matriks pemetaan yang dipakai LatentMAS.

\section{Metodologi Penelitian}

Metodologi yang digunakan adalah studi literatur dan eksperimen terkendali.
Studi literatur dilakukan dengan menelaah publikasi mengenai penalaran laten,
kolaborasi \textit{multi-agent}, dan evaluasi model bahasa, kemudian
menurunkan persamaan tiap metode dari sumber aslinya. Eksperimen dilakukan
dengan menjalankan sistem \textit{multi-agent} pada beberapa tolok ukur dan
satu domain simbolik, dengan seluruh kondisi memakai soal yang sama sehingga
perbandingannya bersifat berpasangan. Analisis dilakukan menggunakan uji
statistik nonparametrik yang sesuai untuk data berpasangan, disertai selang
kepercayaan \textit{bootstrap} dan koreksi terhadap perbandingan berganda.
Seluruh komputasi dilakukan menggunakan bahasa Python.

\section{Sistematika Penulisan}

Penyusunan laporan ini mengacu pada sistematika berikut.

\noindent\textsc{Bab I Pendahuluan}

Bab ini berisi latar belakang, rumusan masalah, pembatasan masalah, tujuan
dan manfaat penelitian, tinjauan pustaka, metodologi penelitian, dan
sistematika penulisan.

\noindent\textsc{Bab II Landasan Teori}

Bab ini menguraikan dasar teori yang digunakan, mencakup probabilitas dan
distribusi yang relevan, teori informasi, inferensi Bayes, aljabar linear dan
representasi vektor, arsitektur \textit{transformer} beserta KV-\textit{cache},
penalaran laten dan ketiga metode pembentukan representasi laten, sistem
\textit{multi-agent}, konsep evaluasi empiris model bahasa, faktor alpha
beserta metrik evaluasinya, dan pengujian statistik yang dipakai.

\noindent\textsc{Bab III Metodologi Penelitian}

Bab ini menetapkan objek penelitian berupa persamaan langkah laten,
menurunkan kelima persamaan yang dibandingkan ke dalam satu kerangka bersama,
menyatakan rancangan dua sumbu eksperimen, hipotesis, data dan tugas evaluasi,
protokol pengambilan sampel berpasangan, metrik, serta prosedur analisis
statistik.

\noindent\textsc{Bab IV Hasil dan Pembahasan}

Bab ini menyajikan hasil eksperimen beserta pembahasannya, mencakup
perbandingan akurasi antar-persamaan, keterbedaan hasil menurut jenis tugas,
biaya komputasi medium laten, fidelitas simbolik keluaran, bukti mekanisme
yang menjelaskan pola hasil, hasil pada domain faktor simbolik, serta batas
berlaku seluruh temuan.

\noindent\textsc{Bab V Kesimpulan dan Saran}

Bab ini memuat kesimpulan yang menjawab rumusan masalah beserta saran bagi
penelitian selanjutnya.
